
{
\begin{center}
    \textbf{{{\Large Appendix}}}
\end{center}
}

\renewcommand\thesection{\Alph{section}}
\renewcommand\thesubsection{\Alph{section}.\arabic{subsection}}


\appendix
\startcontents[sections]
\printcontents[sections]{l}{1}{\setcounter{tocdepth}{2}}



\vspace{20pt}


\section{Reproducibility Statement}\label{appendix:code}


Our model is built upon the publicly available code repository from GLIGEN~\citep{li2023gligen}\footnote{GLIGEN: \url{https://github.com/gligen/GLIGEN}} and \modelscope{}~\citep{wang2023modelscope}\footnote{ModelScopeT2V: \url{https://modelscope.cn/models/damo/text-to-video-synthesis/summary}, \url{https://github.com/ExponentialML/Text-To-Video-Finetuning/tree/main}}. 
Please see \cref{sec:grounded_video_generation}/\cref{appendix:layout2vid} for model architecture details, \cref{sec:experiments}/\cref{sec:single_scene_evaluation_dataset_appendix}/\cref{sec:multi_scene_evaluation_dataset_appendix} for dataset details.
\textbf{We provide code and video samples (\texttt{video\_samples.html}) in the supplementary material.}




\section{Additional Related Works}\label{appendix:related_works}
\paragraph{Text-to-video generation.}
The text-to-video (T2V) generation task is to generate videos from text descriptions.
Early T2V generation models~\citep{li2017video, li2018storygan} used variational autoencoders (VAE)~\citep{kingma2013auto} 
and generative adversarial networks (GAN)~\citep{goodfellow2020generative},
while multimodal language models~\citep{hong2022cogvideo,wu2022nuwa,villegas2022phenaki,maharana2022storydall,ge2022long,wu2021godiva} and denoising diffusion models~\citep{ho2022imagen,singer2022make,blattmann2023align,khachatryan2023text2video,wang2023gen,yin2023nuwa} have become popular for recent works.
Since training a T2V generation model from scratch is computationally expensive,
recent work often leverages pre-trained text-to-image (T2I) generation models such as Stable Diffusion~\citep{rombach2022high} by finetuning them on text-video pairs~\citep{wang2023modelscope, blattmann2023align}.
While this warm-start strategy enables high-resolution video generation, it comes with the limitation of only being able to generate short video clips, as T2I models lack the ability to maintain consistency through long videos. 
Recent works on long video generation \citep{blattmann2023align, yin2023nuwa, villegas2022phenaki, he2023animate} aim at generating long videos of a few minutes.
However, the generated videos often display the continuation or repetitive patterns of a single action (\eg{}, driving a car) instead of transitions and dynamics of multiple changing actions/events (\eg{}, five steps about how to bake a cake).
In this work, we address this problem of multi-scene video generation with a two-stage framework: using an LLM (\eg{}, GPT-4) to generate a structured \videoplan{} (consisting of stepwise scene descriptions, entities and their layouts) and generating videos using \videomodule{}, a layout-guided text-to-video generation model with consistency control.
Our \videomodule{} infuses layout control and multi-scene temporal consistency into a pretrained T2V generation model via data and parameter-efficient training, while preserving its original visual quality.






\section{Video Planning}\label{appendix:prompt_templates}

\subsection{GPT-4 Prompt Templates}
In this section, we provide the prompt templates we give to our video planner (\cref{sec:grounded_video_generation}).
The \videoplan{} is generated in two steps by prompting GPT-4\footnote{We employ the \texttt{gpt-4-0613} version.} with different in-context examples (we use 1 and 5 in-context examples for the first and second steps, respectively).
In the first step (see \cref{fig:steps_prompt}), we ask GPT-4 to expand a single text prompt into a multi-scene \videoplan{}. Each scene comes with a text description, a list of entities, and a background.
In the second step (see \cref{fig:layout_prompt}), we generate a list of bounding boxes for the entities
in each frame based on the list of entities and the scene description. 
In line with VPGen~\citep{cho2023visual}, we utilize the $[x_0, y_0, x_1, y_1]$ format for bounding boxes, where each coordinate is normalized to fall within the range [0,1].
For in-context examples, we present 0.05 as the minimum unit for the bounding box, equivalent to a 20-bin quantization over the [0,1] range. 






\begin{figure}[h]
\begin{center}
    \includegraphics[width=0.99\linewidth]{img/single_text_to_multi_scene_prompt.pdf}
    \caption{Prompt template for the 1st video planning step (scene descriptions and entities/background generation).}
\label{fig:steps_prompt}
\vspace{-4mm}
\end{center}
\end{figure}


\begin{figure}[h]
\begin{center}
    \includegraphics[width=0.99\linewidth]{img/layout_prompt.pdf}
    \caption{Prompt template for the 2nd video planning step (layout generation).}
\label{fig:layout_prompt}
\vspace{-4mm}
\end{center}
\end{figure}


\subsection{API Cost}
Using GPT-4 tokenizer, the average input/output token lengths of each step are 2K/1K for the first step and 6K/1K for the second step.
Using GPT-4, it takes 0.12 USD and 0.24 USD for the inference of the first and second steps, respectively.




\section{\videomodule{}}\label{appendix:layout2vid}


\subsection{Preliminaries: T2V Generation Backbone}\label{appendix:layout2vid_preliminaries}


We implement \videomodule{} by injecting layout control capability into \modelscope{}~\citep{wang2023modelscope}, a public text-to-video generation model based on Stable Diffusion~\citep{rombach2022high}.
\modelscope{} consists of
(1) a CLIP ViT-H/14~\citep{radford2021learning} text encoder,
(2) an autoencoder, and
(3) a diffusion UNet~\citep{ronneberger2015unet,ho2020denoising}.
Given a $T$ frame video $\bm{x}\in\mathbb{R}^{T\times 3 \times H \times W}$ with video caption $\bm{c}$ and frame-wise layouts $\{\bm{e}\}_{i=1}^T$, \modelscope{} first uses an autoencoder
to encode the video into a latent representation $\bm{z}=\mathcal{E}(\bm{x})$.
The diffusion UNet performs denoising steps in the latent space to generate videos, conditioned on the CLIP text encoder representation of video captions.
The UNet comprises a series of spatio-temporal blocks, each containing four modules: spatial convolution, temporal convolution, spatial attention, and temporal attention. Since the original \modelscope{} does not offer control beyond the text input, we build upon the 2D attention module in the spatial attention module to create `Guided 2D Attention', which allows for spatial control using bounding boxes.






\begin{figure}[t]
\centering
    \includegraphics[width=.65\linewidth]{img/network_spatial_temporal_blocks.pdf}
    \vspace{-3mm}
    \caption{
    \textbf{Spatio-temporal blocks within the diffusion UNet of \textbf{\videomodule{}}}.
    The spatio-temporal block comprises four modules: spatial convolution, temporal convolution, spatial attention, and temporal attention. We adopt settings from \modelscope{}, where (N1, N2, N3, N4) are set to (2, 4, 2, 2). We use ``Act'' to represent activation layers.
    }
\label{fig:network}
\end{figure}



\subsection{Layout2Vid Spatio-Temporal Blocks}
\label{appendix:spatio-temporal-blocks}
\cref{fig:network} illustrates the spatio-temporal blocks within the diffusion UNet of \videomodule{}.
The spatio-temporal block comprises four modules: spatial convolution, temporal convolution, spatial attention, and temporal attention.
Following~\cite{li2023gligen}, we build upon the 2D attention to create the guided 2D attention, which enables layout-guided video generation 
with explicit spatial control over a list of entities represented by their bounding boxes, as well as visual and text content.

\subsection{Training and Inference Details}\label{appendix:training_inference_details}

\Cref{table:hyperparameters_gen} contains the model architecture details and training/inference parameter settings for our \videomodule{}.

\paragraph{Training.} The highlight of our \videomodule{} training is that it was conducted solely on image-level data with bounding box annotations.
We trained the MLP layers for grounding tokens and the gated self-attention layers with the same bounding-box annotated data used in GLIGEN~\citep{li2023gligen}, which consists of 0.64M images. 
We train \videomodule{} for 50k steps, which takes only two days with 8 A6000 GPUs (each 48GB memory). 
All the remaining modules in the spatio-temporal block (see main paper Fig. 4) are frozen during training. 
We illustrate the training and inference procedure of \videomodule{} in main paper Fig. 3. 

\paragraph{Inference.} During inference,
we use the Karlo implementation of unCLIP Prior
to the entities to convert the texts into their corresponding CLIP image embeddings, and CLIP text encoder to get their corresponding text embeddings. We use CLIP ViT-L/14 as a backbone during training to be consistent with Karlo. This helps us to preserve the visual consistency of the same object by using the same image embedding across scenes. In addition, we use PLMS~\cite{stablediffusion} (based on PNDM~\cite{liu2022pseudo}) as our default sampler following GLIGEN~\citep{li2023gligen}, since we found no consistent improvements for visual quality and video-text alignment scores when switching to the DDIM~\citep{song2020denoising} sampler used in ModelScopeT2V. The comparison of the PLMS and DDIM samplers is shown in \Cref{table:ablation_samplers}.



\begin{table}[h]
\setlength{\tabcolsep}{4pt}
\centering
\resizebox*{0.72\textwidth}{!}{
    \begin{tabular}{lc@{\hskip 0.75in}lc}
    \toprule
    \textbf{\emph{Architecture}} & & \textbf{\emph{Training}} \\
    $z$-shape & $32\times32\times4$ & \# Train steps & 50k \\
    Channels & 320 & Learning rate & 5e-5\\
    Depth & 2 & Batch size per GPU & 32 \\
    Attention scales & 1,0.5,0.25 & \# GPUs & 8 \\
    Number of heads & 8 & GPU-type & A6000 (48GB) \\
    Head dim & 64 & Loss type & MSE\\
    Channel multiplier & 1.2,4,4 & Optimizer & AdamW\\
    \midrule
    \textbf{\emph{CA Conditioning}} & & \textbf{\emph{Inference}} \\
    Context dimension & 1024 & Sampler & PLMS \\
    Sequence length & 77 & \# denoising steps & 50\\
    \bottomrule
\end{tabular}
}
\caption{Hyperparameters for \videomodule{}.}
\label{table:hyperparameters_gen}
\end{table}






\begin{table}[h]
\begin{center}
    \resizebox{0.5\linewidth}{!}{
    \begin{tabular}{lccc}
        \toprule
        Samplers & FVD () & FID () & CLIPSIM ($\uparrow$) \\
        \midrule
        DDIM  & \textbf{508} & {12.52} & 0.2848  \\
        PLMS  & {550} & \textbf{12.22} & \textbf{0.2860}  \\
        \bottomrule
    \end{tabular}
    }
\end{center}
\caption{
Comparison between PLMS and DDIM sampler on MSR-VTT.}
\label{table:ablation_samplers}
\end{table}














\section{Human-in-the-Loop Video Plan Editing}
\label{appendix:human_in_the_loop}

One useful property of our \method{} is that it is very flexible to edit \videoplan{}s via human-in-the-loop editing.
Given an LLM-generated \videoplan{}, users can generate customized content by adding/deleting/replacing the entities, adding/changing the background, and modifying the bounding box layouts of entities.
\Cref{fig:human_edit} showcases the example videos generated from the text prompt ``A horse running'', and videos generated from the \videoplan{} modified by users (\ie, changing the object sizes and backgrounds).




\begin{figure}[h]
\begin{center}
    \includegraphics[width=.9\linewidth]{img/human_edit.pdf}
    \vspace{-2mm}
    \caption{
        Video generation examples for human-in-the-loop editing.
        Users can modify the video plan (\eg, add/delete objects, change the background and entity layouts, \etc{}) to generate customized video contents.
        Given the same text prompt ``A horse running'', we provide visualizations with a smaller horse and different backgrounds (\ie, ``night street'' and ``grassland'').
    }
\label{fig:human_edit}
\end{center}
\end{figure}




\section{Experiment Setup}\label{sec:experiments_appendix}
We provide additional details on our experiment setups (\cref{sec:experiments} in the main paper)
below.

\subsection{Evaluated Models}
\label{sec:baseline_models_appendix}

We compare our \method{} to 6 popular T2V generation models, 
NUWA~\citep{wu2022nuwa1}, 
CogVideo~\citep{hong2022cogvideo},
VideoLDM~\citep{blattmann2023align},
MagicVideo~\citep{zhou2022magicvideo}, 
Make-A-Video~\citep{singer2022make},
and \modelscope{}~\citep{wang2023modelscope}.
Since NUWA, VideoLDM, MagicVideo, Make-A-Video, and CogVideo (English) are not publicly available,
we primarily compare our \method{} with \modelscope{}, and present comparisons with the other models on the datasets for which their papers have provided results.
\modelscope{} is the closest baseline to our framework among all these models, because our \videomodule{} utilizes its frozen weights and only trains a small set of new parameters to add spatial control and temporal consistency across multiple scenes.

\subsection{Prompts for Single-Scene Video Generation}
\label{sec:single_scene_evaluation_dataset_appendix}


For single-scene video generation,
we conduct experiments with
VPEval Skill-based prompts to evaluate layout control~\citep{cho2023visual},
ActionBench-Direction prompts to assess object dynamics~\citep{wang2023paxion},
and MSR-VTT to cover diverse open-domain scenes~\citep{xu2016msr-vtt}.


\textbf{VPEval Skill-based prompts} 
evaluate different object-centric layout control skills in text-to-image/video generation.
We randomly sample 100 prompts for each of the four skills:
Object (generation of a single object),
Count (generation of a specific number of objects),
Spatial (generation of two objects with a spatial relation; \eg{}, left/right/above/below),
and Scale (generation of two objects with a relative scale relation; \eg{}, bigger/smaller/same).


\textbf{ActionBench-Direction prompts} 
evaluate the action dynamics (object movement directions) in video language models.
We prepare the prompts by sampling video captions from ActionBench-SSV2~\citep{wang2023paxion} and balancing the distribution of movement directions.
Concretely, we select captions from the ActionBench-SSV2 validation split that include phrases like `right to left' or `left to right' (\eg{}, `pushing a glass from left to right'), which are common phrases describing movement directions in the captions.
Then we augment these prompts by switching the directions to each of four directions: `left to right', `right to left', `top to bottom', and `bottom to top' to create 100 prompts for each direction.
We call the resulting 400 prompts as \textit{ActionBench-Direction} prompts. These prompts ensure a balanced distribution of movement directions while maintaining diversity in objects.

\textbf{MSR-VTT} is an open-domain video captioning dataset, which allows us to check if 
our \videomodule{} maintains the original visual quality and text-video alignment performance of the \modelscope{} backbone after integration of the layout/movement control capabilities.
The MSR-VTT test set comprises 2,990 videos, each paired with 20 captions.
Following VideoLDM~\citep{blattmann2023align}, we sample one caption from the 20 available captions for each video and use the 2,990 corresponding generated videos for evaluation. 


\subsection{Prompts for Multi-Scene Video Generation}
\label{sec:multi_scene_evaluation_dataset_appendix}

For multi-scene video generation,
we experiment with two types of input prompts:
(1) a list of sentences describing events -- ActivityNet Captions~\citep{krishna2017dense} and
\corefdata{} prompts based on Pororo-SV~\citep{li2018storygan}
and
(2) a single sentence from which models generate multi-scene videos  -- HiREST~\citep{zala2023hierarchical}.


\textbf{ActivityNet Captions} is a dense-captioning dataset designed for detecting and describing multiple events in videos using natural language.
For the multi-scene video generation task,
we use 165 randomly sampled videos from the validation split
and use the event captions as input for \modelscope{} and our \method{}.
When calculating object consistency, we find the subject of the first event caption via spaCy dependency parser~\citep{spacy2} and check its appearance in multiple scenes.


\textbf{\corefdata{}}
is a new multi-scene text description dataset
that we propose to 
evaluate the consistency of object appearances across multi-scene videos.
We prepare the \corefdata{} prompts by augmenting the Pororo-SV dataset~\citep{li2018storygan,PororoQA2017Kim},
which consists of multi-scene paragraphs from the ``Pororo the Little Penguin'' animated series.
To evaluate the temporal consistency of video generation models trained on real-world videos, we extend its original animation characters (\eg{}, Pororo) to humans and common 
animals and examine their appearance across different scenes.
Concretely, we sample 10 episodes, each consisting of multiple scenes (6.2 scenes per episode on average).
Then, we replace the first appearance of a character with one of the predefined 10 real-world entities (\eg{}, person/dog, etc.) and replace the remaining appearances of the character with pronouns (\eg{}, he/she/it/etc.).
In total, we obtain 100 episodes (=10 episodes $\times$ 10 entities) in \corefdata{}.
In order to generate visually consistent entities,
the multi-scene video generation models would need to address the co-reference of these target entities across different scenes.
We use the final scene descriptions as input for both \modelscope{} and \method{}. When calculating object consistency, we use the selected 
entity as the target object.


\textbf{HiREST} provides step annotations for instructional videos paired with diverse `How to' prompts (\eg{}, a video paired with `how to make butter biscuits' prompt is broken down into a sequence of short video clips of consecutive step-by-step instructions).
For the multi-scene video generation task, we employ 175 prompts from the test splits,
where we only include the prompts with step annotations, to ensure that it is possible to create multi-scene videos from the prompts.
Note that instead of providing a list of scene description sentences like in ActivityNet Captions/\corefdata{}, we only give the single high-level 
`How to' prompt and let the models generate a multi-scene video from it.
In \method{}, our LLM can automatically 
generate the multi-scene \videoplan{} and video from the input prompt.
In contrast, for the \modelscope{} baseline, we help the model understand the different number of scenes to generate by
pre-defining the number of scenes $N$, and independently generate $N$ videos by appending the suffix ``step $n/N$'' to the prompt for $n$-th scene (\eg{}, ``Cook Beet Greens, step 1/10'').
To ensure that our \method{} videos and \modelscope{} videos are equal in length,
we use the same number of scenes generated by our LLM during the planning stage for \modelscope{}. As mentioned in main paper \cref{sec:multi_scene_video_gen}, we also try an ablation with giving \modelscope{} the LLM generated prompts instead of just ``step $n/N$''.









\section{Human Evaluation Details}
\label{appendix:human_eval}

We provide details of our human evaluation and error analysis setups, as well as the error analysis results described in \cref{sec:experiments}.

\paragraph{Human evaluation details.}
We conduct a human evaluation study on the multi-scene videos generated by both our \method{} and \modelscope{} on the \corefdata{} dataset.
Since we know the target entity and its co-reference pronouns in the \corefdata{} prompts,
we can compare the temporal consistency of the target entities across scenes.
We evaluate the human preference between videos from two models in each category of Quality, Text-Video Alignment, and Object Consistency:
\begin{itemize}
    \item \textit{Quality:} it measures how well the video looks visually.
    \item \textit{Text-Video Alignment:} it assesses how accurately the video adheres to the input sentences.
    \item \textit{Object Consistency:} it evaluates how well the target object maintains its visual consistency across scenes.
\end{itemize}
We show 50 videos to ten crowd-annotators from AMT\footnote{Amazon Mechanical Turk: \url{https://www.mturk.com}} to rate each video (28 unique annotators, ten annotators rate each video) and calculate human preferences for each video with average ratings. All videos are randomly shuffled such that annotators do not know which model generated each video.
To ensure high-quality annotations, we require they have an AMT Masters, have completed over 1000 HITs, have a greater than 95\% approval rating, and are from one of the United States, Great Britain, Australia, or Canada (as our task is written in the English language). We pay annotators \$0.10 to evaluate a video (roughly \$12-14/hr).




\paragraph{Step-by-step error analysis details.}
We conduct an error analysis on each step of our single sentence to multi-scene video generation pipeline for HiREST prompts. We analyze the generated multi-scene text descriptions,
layouts, and entity/background consistency groupings to evaluate our video planning stage and the final video to evaluate the video generation stage.
\begin{itemize}
    \item \textit{Multi-Scene Text Descriptions Accuracy}: we measure how well these descriptions depict the intended scene from the original prompt (\eg{}, if the original prompt is ``Make buttermilk biscuits'' the descriptions should describe the biscuit-making process and not the process for pancakes).
    \item \textit{Layout Accuracy}: we measure how well the generated layouts showcase a scene for the given multi-scene text description (\eg{}, the bounding boxes of ingredients should go into a bowl, pan, etc. instead of randomly moving across the scene).
    \item \textit{Entity/Background Consistency Groupings Accuracy}: we measure how well the generated entities and backgrounds are grouped (\eg{}, entities/backgrounds that should look consistent throughout the scenes should be grouped together).
    \item \textit{Final Video Accuracy}: we measure how well the generated video for each scene matches the multi-scene text description (\eg{}, if the multi-scene text description is ``preheating an oven'', the video should accurately reflect this).
\end{itemize}
We ask an expert annotator to rank the generations (multi-scene text description, layouts, etc.) on a Likert scale of 1-5 for 50 prompts/videos. Analyzing the errors at each step enables us to check which parts are reliable and which parts need improvement. As a single prompt/video can contain many scenes, to simplify the process for layout and final video evaluation of a prompt/video, we sub-sample three scene layouts and corresponding scene videos and average their scores to obtain the ``Layout Accuracy'' and ``Final Video Accuracy.''

\paragraph{Step-by-step error analysis results.}
As shown in \Cref{table:error_analysis}, our LLM-guided planning scores high accuracy on all three components (up to 4.52), whereas the biggest score drop happens in the layout-guided video generation (4.52 $\rightarrow$ 3.61).
This suggests that our \method{} could generate even more accurate videos,
once we have access to a stronger T2V backbone than \modelscope{}.


\begin{table}[h]
\begin{center}
\resizebox{0.99\linewidth}{!}{
    \begin{tabular}{ccc c}
        \toprule
        \multicolumn{3}{c}{Stage 1: Video Planning (with GPT-4)} & Stage 2: Video Generation (with \videomodule{}) \\
        \cmidrule(lr){0-2} \cmidrule(lr){4-4} 
        Multi-scene Text Descriptions ($\uparrow$) & Layouts ($\uparrow$) & Entity/Background Consistency Groupings  ($\uparrow$) & Final Video ($\uparrow$) \\
        \midrule
        4.92 & 4.69 & 4.52 & 3.61 \\
        \bottomrule
    \end{tabular}
}
\end{center}
\caption{Step-wise error analysis of \method{} video generation pipeline on HiREST prompts.
We use a Likert scale (1-5) to rate the accuracy of the generated components at each step.
}
\label{table:error_analysis}
\end{table}









\section{Ablation Studies}
\label{sec:ablation_study}



In this section, we provide ablation studies on our design choices,
including using different LLMs for video planning, the number of layout-guided denoising steps,
different embeddings for layout groundings, and layout representation formats.




\subsection{Video Planner LLMs: GPT-4, GPT-3.5-turbo, and LLaMA2}







We conducted an ablation study using the open-sourced LLaMA2-7B and LLaMA2-13B models~\citep{touvron2023llama} as well as GPT-3.5-turbo in the video planning stage on MSR-VTT and ActionBench-Directions. 

\paragraph{Evaluation settings.} In this study, we evaluated LLaMA2 in both fine-tuned and zero-shot settings. For the zero-shot evaluation, we used the same number of in-context examples (\ie, 5 examples) for both GPT-3.5-turbo and GPT-4, but only 3 in-context examples for the LLaMA2 models due to their 4096 token length limits. To collect data for fine-tuning LLaMA2, we randomly sampled 2000 prompts from WebVid-10M~\citep{bain2021frozen} and generated layouts with GPT-4 using the same prompts as described in our main paper. We fine-tuned the LLaMA2 models for 1000 steps with a learning rate of 2e-4 and a cosine scheduler. In addition to the visual quality (FID/FVD) and video-text alignment (CLIPSIM) scores on MSR-VTT and movement direction accuracy score on ActionBench-Directions, we also report the number of samples whose output layouts can be successfully parsed (\eg, \#Samples).
We consider the output valid if it contains 9 frames, and each frame includes at least one object with a bounding box layout. Videos are generated with $\alpha=0.1$ for MSR-VTT and $\alpha=0.2$ for ActionBench-Directions.















\subsection{Entity Grounding Embeddings: Image v.s. Text}
As discussed in \cref{sec:grounded_video_generation} in our main paper,
we compare using different embeddings for entity grounding on 1000 randomly sampled MSR-VTT test prompts.
As shown in \Cref{table:ablation_grounding_token},
CLIP image embedding is more effective than CLIP text embedding,
and using the CLIP image-text joint embedding yields the best results.
Thus, we propose to use the image+text embeddings for the default configuration.


\begin{table}[h]
    \centering
    \resizebox{0.75\linewidth}{!}{
    \begin{tabular}{l ccc c}
        \toprule
        \multirow{2}{*}{\parbox{3.5cm}{Entity Grounding}} & \multicolumn{3}{c}{MSR-VTT} & \corefdata{} \\
        \cmidrule(lr){2-4} \cmidrule(lr){5-5} 
        & FVD () & FID () & CLIPSIM ($\uparrow$) & Consistency (\%) \\
        \midrule
        Image Emb. & {737} & {18.38} & 0.2834 & 42.6  \\
        Text Emb. & {875} & {23.18} & 0.2534 & 36.9 \\
        \rowcolor{lightblue}
        Image+Text Emb. (default) & \textbf{606} & \textbf{14.60} & \textbf{0.2842} & \textbf{42.8} \\
        \bottomrule
    \end{tabular}
    }
    \caption{Ablation of entity grounding embeddings of our \videomodule{} on MSR-VTT and Coref-SV.
    }
    \label{table:ablation_grounding_token}
\vspace{-3pt}
\end{table}
\normalsize 





\subsection{Layout Control: Bounding Box v.s. Center Point}
In \Cref{table:ablation_spatial_control},
we compare different layout representation formats on 1000 randomly sampled MSR-VTT test prompts.
We use image embedding for entity grounding and $\alpha=0.2$ for layout control.
Compared with no layout (`w/o Layout input') or center point-based layouts (without object shape, size, or aspect ratio),
the bounding box based layout guidance gives better visual quality (FVD/FID) and text-video alignment (CLIPSIM).
















\begin{table}[h]
\begin{center}
    \resizebox{0.58\linewidth}{!}{
    \begin{tabular}{lccc}
        \toprule
        Layout representation & FVD () & FID () & CLIPSIM ($\uparrow$) \\
        \midrule
        w/o Layout input  & {639} & {15.28} & 0.2777  \\
        Center point  & {816} & {18.65} & 0.2707  \\
        \rowcolor{lightblue}
        Bounding box (default)  & \textbf{606} & \textbf{14.60} & \textbf{0.2842}  \\
        \bottomrule
    \end{tabular}
    }
\end{center}
\caption{
Ablation of layout representation of our \method{} on MSR-VTT.
}
\label{table:ablation_spatial_control}
\end{table}




\subsection{Trainable Layers: Gated Self-Attention Only v.s. Entire Guided 2D Attention}


\paragraph{Training settings.} In our \videomodule{}, only the MLP layers for grounding tokens and the gated self-attention layers are trained, accounting for 13\% of the total parameters. We conducted an ablation study to unfreeze more parameters from the \modelscope{} backbone. Specifically, all layers in the guided 2D attention module, such as self-attention, gated self-attention, and cross-attention, are made trainable, which constitutes 27\% of the total parameters. Starting from the model trained in our main paper, we further train this variant with more trainable parameters for 50k steps, keeping all other training hyper-parameters constant. 



\begin{table}[h]
\begin{center}
    \resizebox{0.9\linewidth}{!}{
    \begin{tabular}{lccc}
        \toprule
        Trainable Layers & FVD () & FID () & CLIPSIM ($\uparrow$) \\
        \midrule
        Entire Guided 2D Attn (GatedSA + SelfAttn + CrossAttn)  & {1626} & {70.47} & 0.2231  \\
        \rowcolor{lightblue}
         GatedSA only (default)  & \textbf{550} & \textbf{12.22} & \textbf{0.2860}  \\
        \bottomrule
    \end{tabular}
    }
\end{center}
\caption{
Comparison of only updating parameters in gated self-attention (GatedSA) between updating all parameters in guided 2D attention (Entire Guided 2D Attn) of our \method{} on MSR-VTT. 
}
\label{table:ablation_trainable_parameters}
\end{table}





\paragraph{Result analysis.} As shown in \Cref{table:ablation_trainable_parameters}, unfreezing all parameters in the guided 2D attention module results in significantly worse performance compared to training only the additional gated self-attention layers. This decline in performance occurs because the temporal attention and temporal convolution layers, which are inactive during our image-data training, cause misalignment in the spatial and temporal layers during inference when the temporal layers are activated.




\section{Additional Experiments}
\label{appendix:additional_experiments}

\subsection{Comparison on UCF-101}




In addition to the MSR-VTT dataset for open-domain single-scene video generation, we also conduct experiments for the UCF-101 dataset~\citep{soomro2012ucf101}. Since ModelScopeT2V does not report their evaluation results on UCF-101 dataset in their technical report, we use their publicly released checkpoint and run evaluation on the same 2048 randomly sampled test prompts as used in our VideoDirectorGPT evaluation. Compared with ModelScopeT2V, we achieve a significant improvement in FVD (ModelScopeT2V 1093 vs. Ours 748) and competitive performance in the Inception Score (ModelScopeT2V 19.49 vs. Ours 19.42). In addition, we notice that ModelScopeT2V is not good at video generation on this UCF-101 dataset, which is also observed in other recent works~\citep{he2023animate}. We could expect our VideoDirectorGPT to perform better with stronger T2V backbones.



\begin{table}[h]
\centering
\resizebox{0.5\linewidth}{!}{
\begin{tabular}{l cc}
\toprule
\multirow{2}{*}{Method}  & \multicolumn{2}{c}{UCF-101} \\
\cmidrule(lr){2-3} 
& FVD () & IS ($\uparrow$) \\
\midrule
\multicolumn{2}{l}{\textcolor{gray}{Different arch / Training data}} \\
LVDM  &	917 & ${-}$  \\
CogVideo (Chinese)  & 751 & 23.55\\
CogVideo (English) & 701 & 25.27\\
MagicVideo & 699 & ${-}$\\
VideoLDM   & 550 & 33.45	\\
Make-A-Video  & 367  & 33.00\\
\midrule
\multicolumn{3}{l}{\textcolor{gray}{Same video backbone \& Test prompts}} \\
\modelscope{}$^\dagger$  & 1093 & \textbf{19.49} \\
\rowcolor{lightblue}
\method{}  & \textbf{748} & 19.42  \\
\bottomrule
\end{tabular}
}
\caption{ 
Evaluation results on UCF-101.
\modelscope{}$^\dagger$:
Our replication with 2048 randomly sampled test prompts.
}
\label{table:ucf101}
\end{table}




\section{Additional Visualization and Qualitative Examples}
\label{appendix:visualizations}




\subsection{Visualization of Different Layout-Guided Denoising Steps.}\label{appendix:denoising_steps}
\paragraph{Balancing layout control strength with visual quality.}
During video generation,
we use two-stage denoising in \videomodule{} following \citet{li2023gligen},
where we first use layout-guidance with Guided 2D attention for $\alpha*N$ steps,
and use the denoising steps without Guided 2D attention for the remaining $(1-\alpha)*N$ steps, where $N$ is the total number of denoising steps, and $\alpha \in [0,1]$ is the ratio of layout-guidance denoising steps. In our additional analysis (see main paper Table 4), we find that a high $\alpha$ generally increases layout control but could lead to lower visual quality, which is also consistent with the finding in \citet{li2023gligen}.
By default, we use $\alpha=0.1$ and $N=50$ denoising steps.
We also explore using the LLM to determine the $\alpha$ value within the range $[0, 0.3]$ during the \videoplan{} generation
(see main paper \cref{sec:additional_analysis} for details).

\paragraph{Visualization and Analysis.} Here, we provide visualizations of objects with prompts that involve explicit layout movements to illustrate the effect of changing the $\alpha$ hyper-parameter. As can be seen in \cref{fig:visualization_alpha_moveable}, 
a small $\alpha$ value (\eg, 0.1) can provide \textbf{guidance} for object movement, enabling the \videomodule{} to generate videos that do not strictly adhere to bounding box layouts. This approach fosters better creativity and diversity in the generated videos, and enhances visual quality as well as robustness to LLM-produced layouts.
On the other hand, an $\alpha$ value of 0.3 or higher is sufficient to \textbf{control} objects to follow bounding box trajectories. 
In conclusion, for prompts requiring explicit movement control, a larger $\alpha$ value typically results in better performance (as demonstrated by the results from ActionBench-Directions in Table 4 of our main paper). Meanwhile, for prompts that do not require large movements, a smaller 
$\alpha$ value generally leads to improved visual quality (as indicated by the results from MSR-VTT in \Cref{tab:ablation_alpha} in the main paper).



\begin{figure}[t]
\begin{center}
    \includegraphics[width=.88\linewidth]{img/visualization_alpha.pdf}
    \vspace{-2mm}
    \caption{
       Visualization of objects with different strength of layout control (different $\alpha$ values). The prompts are ``A dog moving from right to left'' and ``A car moving from left to right'' respectively. Increasing the $\alpha$ value can make the objects better following bounding box layouts. An $\alpha$ value between 0.1 to 0.3 usually gives videos with best trade-off between visual quality and layout control.
    }
\label{fig:visualization_alpha_moveable}
\end{center}
\end{figure}



\subsection{Additional Qualitative Examples}\label{appendix:additional_qualitative_examples}

\vspace{5pt}


\begin{figure}[t]
\begin{center}
    \includegraphics[width=0.8\linewidth]{img/R1W4.pdf}
\caption{Generated layouts from prompts that require physical understanding: ``A stone thrown into the sky'' (left) and ``A car is approaching from a distance'' (right).}
\label{fig:r1w4}
\end{center}
\end{figure}


\begin{figure}[t]
\begin{center}
    \includegraphics[width=0.99\linewidth]{img/R2W2_3examples.pdf}
\caption{Generated videos from movement prompts with a static object (bottle) and objects that can naturally move (airplan/boat). We use prompts ``A \{bottle/airplane/boat\} moving from left to right.''}
\label{fig:r2w2}
\end{center}
\end{figure}
\vspace{-10pt}




\begin{figure}[h!]
\begin{center}
    \includegraphics[width=0.99\linewidth]{img/R2W3.pdf}
\caption{Multi-scene video generation example from prompt: ``A boy playing in sand and flying kites.''}
\label{fig:multi_scene_boy}
\end{center}
\end{figure}




\paragraph{Prompts requiring understanding of physical world.}
{\cref{fig:r1w4}} shows that our GPT-4 based video planner
can generate object
movements requiring an understanding of the physical world such as gravity and perspective.




\paragraph{Object movement prompts with different types of objects.}
In \cref{fig:r2w2}, we show additional examples for single-scene video generation with prompts involve movements. 
For static objects (\eg{}, a bottle), movement is often depicted via the camera.
For objects that can naturally move (\eg{}, an airplane and a boat), the videos show the movements of objects.




\begin{figure}[h]
\begin{center}
    \includegraphics[width=0.99\linewidth]{img/image_input_example.pdf}
    \vspace{-4mm}
    \caption{
    Video generation examples with \textbf{text-only and image+text inputs}.
    Users can flexibly provide either text-only or image+text descriptions to place custom entities when generating videos with \method{}.
    For both text and image+text based entity grounding examples,
    the identities of the provided entities are well preserved across multiple scenes. 
    }
\label{fig:image_input}
\end{center}
\end{figure}









\paragraph{Generating videos with custom image exemplars.}
In \cref{fig:image_input}, we demonstrate that our Layout2Vid can flexibly take either text-only or
image+text descriptions as input to generate multi-scene videos with good entity consistency.





\paragraph{Multi-scene video generation with human appearance.}
In \cref{fig:multi_scene_boy}, we show an additional example of multi-scene video generation with a human appearance. Our \method{} can well preserve the boy's appearances across scenes.



\paragraph{VPEval Skill-based.}
\cref{fig:vpeval_outputs} shows that LLM-generated video plan successfully guides the Layout2Vid module to accurately place objects in the correct
spatial relations and to generate the correct number of objects.
In contrast, ModelScopeT2V fails to generate
a ‘pizza’ in the first example and overproduces the number of frisbees in the second example.




\begin{figure}[h]
\begin{center}
    \includegraphics[width=.85\linewidth]{img/vpeval_generation_comparison.pdf}
    \caption{
        Video generation examples on \textbf{VPEval Skill-based prompts} for spatial and count skills.
        Our \videoplan{}, with object layouts overlaid, successfully guides the \videomodule{} module to place objects in the correct spatial relations and to depict the correct number of objects,
        whereas \modelscope{} fails to generate a `pizza' in the first example and overproduces the number of frisbees in the second example. 
    }
\label{fig:vpeval_outputs}
\vspace{-2mm}
\end{center}
\end{figure}





\begin{figure}[h]
\begin{center}
    \includegraphics[width=.85\linewidth]{img/actionbench_direction_generation_comparison.pdf}
    \vspace{-2mm}
    \caption{
        Video generation examples on \textbf{ActionBench-Direction prompts}.
        Our \videoplan{}'s object layouts (overlaid) can guide the \videomodule{} module to place and move the `stuffed animal' and `pear' in their correct respective directions,
        whereas the objects in the \modelscope{} videos stay in the same location or move in random directions. 
    }
\label{fig:actionbench_outputs}
\vspace{-2mm}
\end{center}
\end{figure}



\paragraph{ActionBench-direction.}
\cref{fig:actionbench_outputs} shows that
our LLM-generated video plan can guide the Layout2Vid module to place the ‘stuffed animal’ and the ‘pear’
in their correct starting positions and then move them toward the correct end positions, whereas the objects
in the ModelScopeT2V videos stay in the same location or move in random directions.





\begin{figure}[h]
\begin{center}
    \includegraphics[width=.88\linewidth]{img/pororo_generation_comparison.pdf}
    \vspace{-2mm}
    \caption{
        Video generation examples on \textbf{\corefdata{} prompts}.
        Our \videoplan{}'s object layouts (overlaid) can guide the \videomodule{} module to generate the same brown dog and maintain snow across scenes consistently,
        whereas \modelscope{} generates different dogs in different scenes and loses the snow after the first scene. 
    }
\label{fig:pororo_outputs}
\vspace{-2mm}
\end{center}
\end{figure}






\paragraph{\corefdata{}.}
\cref{fig:pororo_outputs} shows that 
our \videoplan{} guide the \videomodule{} module to generate the same dog and maintain snow across scenes consistently,
whereas \modelscope{} generates different dogs in different scenes and loses the snow after the first scene.



\begin{figure}[h]
\begin{center}
    \includegraphics[width=.88\linewidth]{img/hirest_generation_comparison_figure_3.pdf}
    \vspace{-2mm}
    \caption{
        Video generation examples \textbf{HiREST prompts}. Our \method{} generates a detailed \videoplan{} that properly expands the original text prompt, ensures accurate object bounding box locations (overlaid), and maintains consistency of the person across the scenes. 
        \modelscope{} only generates strawberries, without strawberry surprise. 
    }
\label{fig:hirest_outputs_2}
\vspace{-2mm}
\end{center}
\end{figure}







\paragraph{HiREST.}
\cref{fig:hirest_outputs_2} shows that
our LLM can generate step-by-step \videoplan{} from a single prompt, and our \videomodule{} can generate consistent videos following the plan.
Our \method{} breaks down the process and generates a complete video showing how to make a strawberry surprise (a type of dessert consisting of vanilla ice cream and strawberries). \modelscope{} repeatedly generates strawberries.











\section{Limitations}
\label{sec:limitations}


While our framework can be beneficial for numerous applications (\eg{}, user-controlled/human-in-the-loop video generation/manipulation and data augmentation), akin to other video generation frameworks, it can also be utilized for potentially harmful purposes (\eg{}, creating false information or misleading videos), and thus should be used with caution in real-world applications (with human supervision, \eg{}, as described in \Cref{appendix:human_in_the_loop} --- human-in-the-loop video plan editing).
Also, generating a \videoplan{} using the strongest LLM APIs can be costly,
similar to other recent LLM-based frameworks.
We hope that advances in quantization/distillation and open-source models will continue to lower the inference cost of LLMs.
Lastly, our video generation module (\videomodule{}) is based on the pretrained weights of a specific T2V generation backbone. Therefore, we face similar limitations to their model, including deviations related to the distribution of training datasets, imperfect generation quality, and only understanding the English corpus.
\vspace{20pt}






